\documentclass[11pt,a4paper]{report}
\usepackage[utf8]{inputenc}
\usepackage[T1]{fontenc}
\usepackage{lmodern}
\usepackage[french]{babel}
\usepackage{geometry}
\geometry{margin=2.5cm,headheight=14pt}
\usepackage{parskip}
\usepackage{fancyhdr}
\usepackage{hyperref}
\hypersetup{hidelinks,pdftitle={$title$},pdfauthor={$author$}}
\providecommand{\tightlist}{\setlength{\itemsep}{0pt}\setlength{\parskip}{0pt}}
\providecommand{\releaseversion}{Édition de travail}
\setlength{\emergencystretch}{3em}
\pagestyle{fancy}
\fancyhf{}
\fancyhead[L]{\small Léon Peillon}
\fancyhead[R]{\small La famille Berloty}
\fancyfoot[C]{\thepage}
\begin{document}
\begin{titlepage}
\centering
\vspace*{1.8cm}
{\Large Léon Peillon\par}
\vspace{1.8cm}
{\LARGE\bfseries Des débuts modestes\par}
\vspace{0.6cm}
{\Large de Claude Denis Berloty\par et Scholastique Bonaventure Poulat\par}
\vspace{0.8cm}
{\Large à la fortune immobilière\par laissée par leur fils F.F. Berloty en 1871\par}
\vfill
{\large Février 1988\par}
\vspace{1cm}
{\small Transcription des onze pages du tapuscrit familial\par}
{\small \releaseversion\par}
\end{titlepage}
\tableofcontents
\clearpage
$body$
\end{document}
